% LaTeX Curriculum Vitae
% Author: Xovee Xu (https://xovee.cn)
% Last Update: July 28, 2020
% Latest Version (consider give me a star!): https://github.com/Xovee/latex-cv
% Overleaf Link: https://www.overleaf.com/latex/templates/xovees-cv-template/rrsmqwhbygcf
% Current Version: v0.95
% LICENSE: MIT


\documentclass{article}

\usepackage[utf8]{inputenc}
\usepackage[full]{textcomp}
\usepackage[lf]{ebgaramond}
% \usepackage[scaled,swashQ]{garamondx}
\usepackage[T1]{fontenc}

\usepackage{enumitem}
\usepackage[a4paper,left=.9in, right=.9in, top=.85in, bottom=.85in]{geometry}
\usepackage{url}
\usepackage[dvipsnames]{xcolor}
\usepackage{amsmath}
\usepackage{titlesec}

% package settings
\usepackage[
    hidelinks,
    pdfnewwindow=true,
    pdfauthor={Shih-Yang Lin},
    pdftitle={Curriculum Vitae of Shih-Yang Lin},
]{hyperref}

\pagestyle{headings}
\markright{\textbf{Shih-Yang Lin}}

\setlength\parindent{2em}
\setlength{\parskip}{2pt}

\thispagestyle{empty}

% section headings: bold caps + thin rule beneath, tightened spacing
\titleformat{\section}{\rmfamily\bfseries\large}{}{0pt}{}[{\titlerule[0.4pt]}]
\titlespacing{\section}{0pt}{10pt}{6pt}
\titlespacing{\subsection}{0pt}{8pt}{4pt}

\newcommand{\cvsection}[1]{\section*{\MakeUppercase{#1}}}
\newcommand{\cvsubsection}[1]{\subsection*{\rmfamily\hspace{1.6em}#1}}

% shared hanging-indent list style used by every entry list in this CV
\newlist{cvlist}{itemize}{1}
\setlist[cvlist]{leftmargin=2em, topsep=2pt, itemsep=3pt, label={}}

% begin
\begin{document}

% name and header
\begin{center}
    {\LARGE\rmfamily\bfseries Shih-Yang Lin}\\[4pt]
    {\normalsize\rmfamily\itshape Assistant Professor, Department of Economics, National Dong Hwa University}\\[8pt]
    {\small\rmfamily No.\,1, Sec.\,2, Da Hsueh Rd., Shoufeng, Hualien 974301, Taiwan}\\[2pt]
    {\small\rmfamily +886-3-890-5537 \quad$\cdot$\quad \href{mailto:linsy@gms.ndhu.edu.tw}{\texttt{linsy@gms.ndhu.edu.tw}} \quad$\cdot$\quad \href{https://orcid.org/0000-0003-0864-0187}{\texttt{ORCID: 0000-0003-0864-0187}}}
\end{center}
\vspace{6pt}
\noindent\rule{\linewidth}{0.8pt}
\vspace{4pt}

\renewcommand{\arraystretch}{1.25}




% Summary / Statement
%\cvsection{Summary / Statement}
%\indent

%\hangindent=2emI have rich experience and strong skills in ... I want ...


% Education
\cvsection{Education}
\indent
\textbf{National Taiwan University (NTU)}, Taipei, Taiwan
\begin{cvlist}
    \item Ph.D., Economics, 2025
    \item M.A., Economics, 2016
\end{cvlist}

\textbf{National Sun Yat-Sen University (NSYSU)}, Kaohsiung, Taiwan
\begin{cvlist}
    \item B.A., Political Economy, 2008
\end{cvlist}


% research interests
\cvsection{Research Interests}
\indent

Macroeconomics, Financial Economics, Econometrics, International Finance


% below is another kind of style
%\cvsection{Education (Another Style)}
%\indent 

%Ph.D., Computer Science, University of Electronic Science and Technology of China, 2021 to 2025 (expc.)

%M.S., Software Engineering, University of Electronic Science and Technology of China, 2021

%B.S., Software Engineering, University of Electronic Science and Technology of China, 2018


% academic employment
\cvsection{Academic Employment}
\indent
\textbf{National Dong Hwa University}, Hualien, Taiwan
\begin{cvlist}
    \item Assistant Professor, Department of Economics, 2025--present
\end{cvlist}

\textbf{Academia Sinica}, Taipei, Taiwan
\begin{cvlist}
    \item Research Assistant, 2016--2018
\end{cvlist}


%\cvsection{Academic Employment (Another Style)}
%\indent

%Dean, Dept. of Computer Science, Standard University, CA, USA, 2080 -

%Professor, Standard University, CA, USA, 2077 -

%Associated Professor, Standard University, CA, USA, 2066 - 2077

%Assistant Professor, Standard University, CA, USA, 2060 - 2066

%Postdoc, University of Electronic Science and Technology of China, Chengdu, China, 2059 - 2060

% distinction, award, honor, and fellowship
\cvsection{Grants and Awards}

\begin{cvlist}
    \item Best Doctoral Dissertation Award (Sho-Chieh Tsiang Theoretical Economics), Taiwan Economic Association, 2025
    \item Doctoral Dissertation Fellowship of the Humanities and Social Sciences, National Science and Technology Council, 2023
    \item Asia Bank Scholarship, National Taiwan University, 2022
    \item Liu Ta-Chung Scholarship, National Taiwan University, 2020
    \item Harold H.C. Han Scholarship, National Taiwan University, 2019
    \item Chu, Cyrus C. Y. and Chen, Tain-Jy Scholarship, National Taiwan University, 2018
\end{cvlist}

% publication
\section*{\MakeUppercase{Publication}}

\cvsubsection{Refereed Journal Articles}

\begin{enumerate}
    \item \textbf{“U.S. Inflation Inequality between 2010 and 2023.”} Coauthored with YiLi Chien and Yi-Chan Tsai. 2026. \textit{Federal Reserve Bank of St. Louis Review}, 108: (2), pp. 1-22,\\ doi:\href{https://doi.org/10.20955/r.2026.02}{10.20955/r.2026.02}
    \item \textbf{“Financial Frictions and Uncertainty Shocks.”} Coauthored with Yi-Chan Tsai and Po-Yuan Wang. 2025. \textit{Macroeconomic Dynamics}, 29: e143,\\ doi:\href{https://doi.org/10.1017/S1365100525100497}{10.1017/S1365100525100497}
    \item \textbf{“Estimating the Willingness to Pay for Voting when Absentee Voting is not Allowed.”} Coauthored with C.Y. Cyrus Chu and Wen-Jen Tsay. 2021. \textit{Social Science Quarterly}, 102: (4), pp. 1380-1393,\\ doi:\href{https://doi.org/10.1111/ssqu.13014}{10.1111/ssqu.13014}
\end{enumerate}


\cvsubsection{Non-refereed Journal Articles}

\begin{enumerate}
    \item \textbf{“Constructing the DSGE-VAR Model — Taiwan's Medium-to-Long-Term Economic Growth Rate Forecast.”} Coauthored with Yi-Chan Tsai, Ruey Yau, and Kuo-Hsuan Chin, 2022. \textit{Quarterly Bulletin, Central Bank of the Republic of China (Taiwan)}, 44: (2), pp. 27-70.
\end{enumerate}

\section*{\MakeUppercase{Working Paper}}

\begin{enumerate}
    \item \textbf{“The Labor Supply of Married Women and Spousal Tax Deductions in Japan.”} Coauthored with Minchung Hsu, Pei-Ju Liao, and Yi-Chan Tsai.
    \begin{itemize}
        \item[] Population aging results in the problem of labor shortage. The labor force of Japanese females presents a solution to the Japanese government. Despite measures to improve the female labor force participation rate, encouraging women to work more or even shift from part-time jobs to full-time jobs was considered by the government. This study employs a standard incomplete markets model (or the Bewley-Huggett-Aiyagari style model) with heterogeneous households to examine the impact of two tax reforms (2004 and 2018) and a counterfactual reform abolishing spousal deduction on women's employment. Findings indicate that reducing spousal deduction after the 2004 tax reform boosts labor participation due to negative wealth effects, while the 2018 income threshold extensions result in a shift from full-time to part-time employment. Conversely, completely abolishing spousal deduction encourages a shift from part-time to full-time employment.
    \end{itemize}
    \item \textbf{“A Two-Stage Pseudo-Maximum Likelihood Method for Poisson Models with Two Binary Endogenous Explanatory Variables.”} Coauthored with Chingmu Chen, Shin-Kun Peng, and Wen-Jen Tsay.
    \begin{itemize}
        \item[] This paper introduces a novel estimator for the Poisson model with dual binary endogenous explanatory variables, addressing the need to quantify the effects of preferential economic integration agreements. We develop a two-stage pseudo-maximum likelihood estimator and derive analytical gradient matrices as well as Hessian matrices. These exact formulas improve computational speed significantly. Applying our approach to trade data, we find that preferential trade agreements have a positive impact on bilateral trade flows, while bilateral investment treaties have a negative effect. Additionally, the positive interaction term between these two treaties suggests that preferential trade agreements promote horizontal foreign direct investment.
    \end{itemize}
    \item \textbf{“Robust Inference for Interrupted Time Series Models with Serial Dependent Disturbances.”} Coauthored with Wen-Jen Tsay and Hsin-Chieh Wong.
    \begin{itemize}
        \item[] This paper considers inference problems related to the widely known interrupted time series (ITS) model, which is used as a quasi-experimental approach for evaluating policy intervention effects over time. Under a single-treatment-group design with serially correlated errors, we are the first to establish the asymptotic distribution of the ordinary least squares (OLS) estimators, which involves Brownian motion functionals and reveals a superconsistency property for the trending variables when the estimation error of the OLS estimator relative to the true coefficients is properly rescaled. Interestingly, the asymptotic covariance matrix of the rescaled estimator is proved to be invertible, even though the convergence rates of the OLS estimator for the level and trending variables differ, and the trending variables have different observation lengths. We further develop a studentized sieve bootstrap procedure and establish its asymptotic validity, providing an implementable inference method for small samples. In simulation experiments, the proposed bootstrap-based inference performs well, particularly under strong serial dependence.\\ doi:\href{https://dx.doi.org/10.2139/ssrn.5486489}{10.2139/ssrn.5486489}
    \end{itemize}
    \item \textbf{“Estimating Count Data Models with $k$ Endogenous Switching Variables.”}
    \begin{itemize}
        \item[] This paper proposes a novel approach to address the estimation and inference problems concerning the count data models featuring $k$ binary endogenous explanatory variables. The proposed estimator fills the gap in various economic domains where discussions involve endogenous dummy variables (treatment effects) and sample selection. Under the correct model specification the full information maximum likelihood approach to estimation is introduced. Additionally, we develop both a one-stage quasi-maximum likelihood estimator and a two-stage quasi-maximum likelihood estimator to offer relatively robust alternatives.
    \end{itemize}
    \item \textbf{“Empirical Relevance of Collateral-based and Cash Flow-based Credit Constraints in the U.S. Economy.”}
    \begin{itemize}
        \item[] This study explores the significance of cash flow-based borrowing constraints, testing their empirical relevance across various settings. It investigates whether these constraints can explain key features of U.S. economic data, such as the lead-lag relationship and volatility of residential and nonresidential investments. The study also compares cash flow-based constraints with asset-based counterparts to determine which is more empirically relevant. Finally, it examines the primary driving force behind the Great Recession, focusing on the impacts of deteriorating credit conditions on households and firms. The aim is to provide a clearer understanding of the main factors responsible for the downturn, informing future policy decisions.
    \end{itemize}
\end{enumerate}



% talks
\cvsection{Conference Presentation}

\begin{cvlist}
    \item \textbf{National Taipei University Interschool Academic Conference --- 2026}, Taipei, Taiwan, January 19, 2026. \\
    \hspace*{1.2em}\textit{“A Two-Stage Pseudo-Maximum Likelihood Method for Poisson Models with Two Binary Endogenous Explanatory Variables: An Application to Trade and Investment Agreements.”}
    \item \textbf{Taiwan Economic Association Annual Conference --- 2025}, Taipei, Taiwan, November 29, 2025. \\
    \hspace*{1.2em}\textit{“Asymptotic Distribution of OLS Estimator for Interrupted Time Series Model.”}
    \item \textbf{Macroeconometric Modelling Workshop --- 2024}, Taipei, Taiwan, December 11--12, 2024. \\
    \hspace*{1.2em}\textit{“The Labor Supply of Married Women and Spousal Tax Deductions in Japan.”}
    \item \textbf{Macroeconometric Modelling Workshop --- 2023}, Taipei, Taiwan, November 27--28, 2023. \\
    \hspace*{1.2em}\textit{“Financial Frictions and Uncertainty Shocks.”}
    \item \textbf{Taiwan Economic Association Conference --- 2016}. \\
    \hspace*{1.2em}\textit{“Estimating the Willingness to Pay for the Civic Value When Absentee Voting Is Not Allowed.”}
\end{cvlist}

\cvsubsection{Discussant}

\begin{cvlist}
    \item Discussant, “Overnight Momentum and Intraday Reversal: Evidence from Taiwan Listed Companies,” National Taipei University Interschool Academic Conference, 2026.
\end{cvlist}


% teaching experience
\cvsection{Teaching Experience}

\cvsubsection{National Dong Hwa University -- Instructor}

\begin{cvlist}[leftmargin=0pt]
    \item \textbf{Macrofinance}\hfill Fall 2026
    \item \textbf{Macroeconomics (I)}\hfill Fall 2026
    \item \textbf{Econometric Analysis (II)}\hfill Spring 2026
    \item \textbf{Introduction to Big Data Analytics}\hfill Fall 2025
    \item \textbf{Development Economics}\hfill Fall 2025
    \item \textbf{Financial Markets}\hfill Spring 2025--2026
    \item \textbf{Financial Economics (I)}\hfill Spring 2025
\end{cvlist}

\cvsubsection{National Taiwan University -- Teaching Assistant}

\begin{cvlist}[leftmargin=0pt]
    \item \textbf{Macroeconomic Theory (I)}\hfill Fall 2019--2022
    \item \textbf{Macroeconomic Theory (II)}\hfill Spring 2020--2021, 2024
    \item \textbf{Theory of Income and Employment (III)}\hfill Fall 2022
    \item \textbf{Economics: Empirical Study and Forecasting (I \& II)}\hfill 2019--2022
    \item \textbf{Basic Macroeconomics}\hfill Summer 2022
    \item \textbf{Basic Quantitative Method}\hfill Summer 2020--2022
    \item \textbf{Principle of Economics (I \& II)}\hfill 2019--2020
\end{cvlist}


% industrial experience
% \section*{Industrial Experience}
% \indent

% \textbf{Strawberry Inc.}, Beijing, China

% \hspace{2em}Co-funder, 2077 - 

% \textbf{Delicious Mango Holdings Ltd.}, London, UK

% \hspace{2em}\LaTeX~Engineer Intern, 2076 - 2077

% \textbf{CV Agency Co. Ltd.}, Hawaii, USA

% \hspace{2em}CV Broker, 2075 - 2076
% \begin{itemize}[leftmargin=5em, topsep=0pt, itemsep=0pt]
%     \item Helped a three-year old child create her first CV
%     \item Hilight 2 ...
%     \item Hilight 3 ...
% \end{itemize}




% language
%\cvsection{Language}
%\indent

%Native Chinese (Mandarin \& Northern Shaanxi Dialect)

%Fluent English




% reference 
\cvsection{Reference}
\indent

\begin{tabular}{lrr}
% Referee 1
\begin{minipage}[t]{4.8cm}
Yi-Chan Tsai                    \\
Department of Economics         \\
National Taiwan University      \\
No. 1, Sec. 4, Roosevelt Road,  \\
Taipei 10617, Taiwan            \\
+886-2-3366-8314                \\
\href{mailto:yichantsai@ntu.edu.tw}{yichantsai\textrm{@}ntu.edu.tw}
\end{minipage}
&
% Referee 2
\begin{minipage}[t]{4.8cm}
Wen-Jen Tsay                    \\
Institute of Economics          \\
Academia Sinica                 \\
No. 128, Sec. 2, Academia Road, \\
Taipei 11529, Taiwan            \\
+886-2-2782-2791\#296           \\
\href{wtsay@econ.sinica.edu.tw}{wtsay\textrm{@}econ.sinica.edu.tw}
\end{minipage}
&
\begin{minipage}[t]{4.8cm}
Pei-Ju  Liao                    \\
Department of Economics         \\
National Taiwan University      \\
No. 1, Sec. 4, Roosevelt Road,  \\
Taipei 10617, Taiwan            \\
+886-2-3366-8371                \\
\href{mailto:pjliao@ntu.edu.tw}{pjliao\textrm{@}ntu.edu.tw}
\end{minipage}
\end{tabular}


%\vfill

%\section*{\hfill\color{OliveGreen}\today}

\end{document}
